% ════════════════════════════════════════════════════════════════
% 1.1 - CONCEITOS DE DERIVADA
% ════════════════════════════════════════════════════════════════

\chapter{Cálculo Diferencial em R}

\section{Conceitos de Derivada}

% ────────────────────────────────────────────────────────
\subsection{Enquadramento Teórico}

\begin{notebox}[title={Intuição Essencial}]
  A derivada mede a \textbf{velocidade instantânea de mudança}. Se uma função descreve uma quantidade que varia (como a posição, a temperatura ou o custo financeiro), a derivada indica-nos quão rápido essa quantidade está a mudar num exato instante.
\end{notebox}

\begin{defbox}[title={Derivada de uma Função num Ponto}]
  Seja $f:\mathbb{R}\to\mathbb{R}$. Diz-se que $f$ é \textbf{derivável} (ou diferenciável) num ponto $x_0$ se o seguinte limite existir e for um número finito:
  \[
    f'(x_0) = \lim_{h\to 0} \frac{f(x_0+h)-f(x_0)}{h}
  \]
  Geometricamente, este valor representa a \textbf{inclinação (ou declive) da reta tangente} à curva representativa de $f$ no ponto de abcissa $x_0$. O termo interno $\frac{f(x_0+h)-f(x_0)}{h}$ é a \emph{taxa média de variação} (ou quociente incremental).
\end{defbox}

\begin{defbox}[title={Regras Fundamentais de Derivação}]
  Sejam $f$ e $g$ funções deriváveis e $c, n \in \mathbb{R}$. As operações com derivadas respeitam as seguintes regras operatórias:
  \begin{itemize}
      \item \textbf{Constante:} $(c)' = 0$
      \item \textbf{Potência:} $(x^n)' = n x^{n-1}$
      \item \textbf{Soma e Diferença:} $(f \pm g)' = f' \pm g'$
      \item \textbf{Multiplicação por Escalar:} $(cf)' = c f'$
      \item \textbf{Produto:} $(f \cdot g)' = f'g + f g'$
      \item \textbf{Quociente:} $\displaystyle \left(\frac{f}{g}\right)' = \frac{f'g - f g'}{g^2}$ \quad (para $g \neq 0$)
      \item \textbf{Regra da Cadeia (Funções Compostas):} $(f \circ g)'(x) = f'(g(x)) \cdot g'(x)$
  \end{itemize}
\end{defbox}

\begin{warnbox}[title={Erros Comuns na Derivação}]
  \begin{itemize}
      \item \textbf{Confundir o quociente incremental com a derivada:} A expressão $\frac{f(x+h)-f(x)}{h}$ dá-nos apenas a velocidade \emph{média} entre dois pontos. A derivada (velocidade \emph{instantânea}) exige a aplicação obrigatória do limite $h \to 0$.
      \item \textbf{Achar que "Contínua $\implies$ Derivável":} Esta implicação é \textbf{falsa}. Toda a função derivável é garantidamente contínua, mas \emph{nem toda} função contínua é derivável. O contraexemplo clássico é $f(x) = |x|$, que não tem derivada na origem devido a uma "quebra" abrupta na geometria.
  \end{itemize}
\end{warnbox}

% ────────────────────────────────────────────────────────
\subsection{Exemplos Resolvidos}

\begin{exbox}{1 --- Derivada de uma Função Potência}
  Calcular a derivada da função $f(x)=x^3$ e interpretar o resultado.
\end{exbox}
\begin{solbox}
  Aplicando diretamente a regra da potência $(x^n)' = n x^{n-1}$ com $n=3$:
  \[
    f'(x) = 3x^{3-1} = \mathbf{3x^2}
  \]
  \textbf{Interpretação:} A inclinação da reta tangente à curva num dado ponto $x$ é dada por $3x^2$. Como $x^2 \ge 0$, a função está sempre a crescer (ou com inclinação nula na origem, onde $x=0$). À medida que o valor de $x$ se afasta do centro, a inclinação aumenta exponencialmente, tornando a curva cada vez mais íngreme.
\end{solbox}

\begin{exbox}{2 --- Falha na Derivabilidade (O "Bico")}
  A função $f(x)=|x|$ não é derivável em $x=0$. Explique o porquê de forma intuitiva.
\end{exbox}
\begin{solbox}
  A função módulo forma um "bico" (ponto anguloso) exatamente em $x=0$.
  Se analisarmos as inclinações das semirretas laterais:
  \begin{itemize}
      \item O lado esquerdo da curva tem inclinação constante de $-1$.
      \item O lado direito da curva tem inclinação constante de $+1$.
  \end{itemize}
  Como as inclinações laterais não coincidem ($-1 \neq 1$), o limite que define a derivada não existe nesse ponto. Na prática, não conseguimos equilibrar uma \emph{única} reta tangente na ponta daquele bico.
\end{solbox}

% ────────────────────────────────────────────────────────
\subsection{Exercícios Práticos}

\begin{exercisebox}{1 --- Aplicação Múltipla de Regras}
  Calcule a função derivada $f'(x)$ para o polinómio $f(x)=5x^4 - 3x + 7$.
\end{exercisebox}
\begin{solbox}
  Utilizando linearmente as regras da soma, da multiplicação por escalar, da potência e da constante:
  \begin{align*}
      f'(x) &= (5x^4)' - (3x)' + (7)' \\
            &= 5(4x^3) - 3(1) + 0 \\
            &= 20x^3 - 3
  \end{align*}
  \textbf{Solução Final:} $\boxed{f'(x) = 20x^3 - 3}$.
\end{solbox}

\begin{exercisebox}{2 --- Análise Analítica de Derivabilidade}
  Determine se a função $f(x)=|x-2|$ é derivável em $x=2$. Justifique o seu raciocínio.
\end{exercisebox}
\begin{solbox}
  Vamos calcular os limites laterais do quociente incremental no ponto $x_0 = 2$.
  Sabemos que $f(2) = |2-2| = 0$.
  
  \textbf{Limite à direita ($h \to 0^+$):}
  \[
    \lim_{h \to 0^+} \frac{f(2+h) - f(2)}{h} = \lim_{h \to 0^+} \frac{|2+h-2|}{h} = \lim_{h \to 0^+} \frac{|h|}{h} = \lim_{h \to 0^+} \frac{h}{h} = \mathbf{1}
  \]
  
  \textbf{Limite à esquerda ($h \to 0^-$):}
  \[
    \lim_{h \to 0^-} \frac{f(2+h) - f(2)}{h} = \lim_{h \to 0^-} \frac{|h|}{h} = \lim_{h \to 0^-} \frac{-h}{h} = \mathbf{-1}
  \]
  
  Como os limites laterais devolvem valores distintos ($1 \neq -1$), o limite global em $x=2$ não existe. \\
  \textbf{Conclusão:} A função \textbf{não é derivável} em $x=2$.
\end{solbox}

\begin{exercisebox}{3 --- Prova pela Definição Formal}
  Mostre, recorrendo estritamente à definição formal de limite, que para $f(x) = x^2$, a derivada é efetivamente $(x^2)' = 2x$.
\end{exercisebox}
\begin{solbox}
  Pela definição, a derivada é dada por:
  \[
    f'(x) = \lim_{h\to 0} \frac{f(x+h)-f(x)}{h}
  \]
  Substituindo a função pela expressão $f(x) = x^2$:
  \begin{align*}
      f'(x) &= \lim_{h\to 0} \frac{(x+h)^2 - x^2}{h} \\
            &= \lim_{h\to 0} \frac{x^2 + 2xh + h^2 - x^2}{h} \\
            &= \lim_{h\to 0} \frac{2xh + h^2}{h}
  \end{align*}
  Fatorizando o $h$ no numerador (visto que $h \neq 0$ durante a aproximação no limite):
  \[
    f'(x) = \lim_{h\to 0} \frac{h(2x + h)}{h} = \lim_{h\to 0} (2x + h)
  \]
  Aplicando a substituição direta do limite quando $h \to 0$, obtemos finalmente:
  \[
    f'(x) = 2x + 0 = \mathbf{2x} \quad \text{\checkmark}
  \]
\end{solbox}

% ────────────────────────────────────────────────────────
\subsection{Questões Propostas para Defesa Oral}

\begin{oralbox}
  \textbf{Q1.} Como explicaria o conceito prático de derivada a alguém do ensino básico? (Sugestão: recorra à analogia de um carro em viagem a medir a velocidade no conta-quilómetros).
\end{oralbox}

\begin{oralbox}
  \textbf{Q2.} Qual a diferença rigorosa entre \emph{taxa média de variação} e \emph{taxa instantânea de variação}? Relacione ambos os conceitos geométricos (reta secante \textit{vs.} reta tangente).
\end{oralbox}

\begin{oralbox}
  \textbf{Q3.} Em que pontos exatos a função modular típica, como $|x|$, falha a derivabilidade e porquê? Explique porque é que a \emph{continuidade} não é suficiente para garantir a derivação.
\end{oralbox}

\begin{oralbox}
  \textbf{Q4.} Explique geometricamente ou através de taxas de variação interdependentes a intuição por trás da Regra da Cadeia. Como se comportam os declives numa função composta?
\end{oralbox}